\subsection{The map-reduce pattern}

\begin{frame}[t,fragile]{The map-reduce pattern}
  \begin{itemize}
    \item Map-reduce is a common pattern for parallel algorithms.
    \item It consists of two steps:
      \begin{itemize}
        \item \textmark{Map}: apply a function to each element of a sequence (or pairs of elements).
        \item \textmark{Reduce}: combine the results of the mapping using an associative and commutative operation.
      \end{itemize}

    \mode<presentation>{\vfill\pause}
    \item \textmark{Variants}:
      \begin{itemize}
        \item \textgood{Unary map-reduce}: map a unary function and reduce with a binary operation.
\begin{lstlisting}
T transform_reduce(f, l, init, binary_comb_op, unary_map_op);
T transform_reduce(ex, f, l, init, binary_comb_op, unary_map_op);
\end{lstlisting}

        \mode<presentation>{\vfill\pause}
        \item \textgood{Binary map-reduce}: map a binary function on pairs of elements and reduce with a binary operation.
\begin{lstlisting}
T transform_reduce(f1, l1, f2, init);
T transform_reduce(ex, f1, l1, f2, init);
T transform_reduce(f1, l1, f2, init, binary_comb_op, binary_map_op);
T transform_reduce(ex, f1, l1, f2, init, binary_comb_op, binary_map_op);
\end{lstlisting}
      \end{itemize}
  \end{itemize}
\end{frame}

\begin{frame}[t,fragile]{Adding squares of numbers}

\begin{block}{Map-Reduce}
\begin{lstlisting}
auto add_squares_par(std::vector<int> const & vec) {
    return std::transform_reduce(std::execution::par, vec.begin(), vec.end(), 0, 
    [](int a, int b) {
      return a + b;
    },
    [](int x) {
      return x*x;
    });
}
\end{lstlisting}
\end{block}

\begin{itemize}
  \item \textgood{Solution}: The new combination operation is now associative and commutative, so the result is deterministic and can be parallelized; the map operation should also be pure.
\end{itemize}

\end{frame}